% Numerical values below are taken from the validated 100-background and
% independently labelled continuation ledgers in the v4.9.1 bundle.

\section{Background-model validation with full-refit pseudoexperiments}
\label{sec:toys}

\subsection{Validation question and scope}

The validation begins with a concrete local question.  If a smooth, data-like spectrum
contains no resonance, does the complete GPR sideband fit return a spurious signal?
If a known signal is added, does the same procedure recover it with an uncertainty of
the right scale?  The pseudoexperiments exercise the count-space prediction,
the mass-dependent blind and training masks, the full-range Gaussian signal
template, the matched-reference injection convention, and a fresh GP optimization
for every pseudo-spectrum.  The statistical definitions are given in
Section~\ref{sec:matched-reference-diagnostics}; this section presents the physical
motivation, the generating spectra, the ensembles, the numerical safeguards, and
the closure evidence.

\begin{figure}[htbp]
\centering
\graphicorplaceholder{0.98\linewidth}{guard_band_figs/hps_gpr_bump_hunt_flowchart.png}
\caption{Analysis and validation flow.  In the data path, sidebands around each mass
hypothesis train the GP and the blinded bins enter the signal-plus-background
likelihood.  In the validation path, a smooth source-fitted mean is Poisson
fluctuated, a known signal may be added, and the complete GP fit and extraction are
repeated.  The zero-injection ensemble tests spurious signal; nonzero injections test
signal recovery and uncertainty calibration.  These are extraction-closure tests,
not confidence-limit coverage or a scan-wide significance calibration.}
\label{fig:v491-hps-gpr-flowchart}
\end{figure}

Three distinctions are important throughout this section.  First, a generating
function is a smooth stress truth fitted to an HPS source spectrum; it is not claimed
to be the unique physical parameterization of the trident continuum.  Second,
closure is conditional on that specified truth and on the analysis card used for the
extraction.  It therefore tests the behavior of the procedure for data-like spectra
without measuring a bias in the observed data.  Third, the four exposure lanes are
not interchangeable.  The $1\%\times10$ and $1\%\times100$ spectra derive from the
native 1\% source, while native 10\% and $10\%\times10$ derive from an independently
selected 10\% source.  Cross-source differences include source and selection shape
as well as normalization.

\subsection{Closure observables and decision rules}

At zero injection, the signed pull
$p=\widehat A/\widehat\sigma_A$ is a spurious-signal statistic.  The ensemble mean
tests whether the extraction is centered and the sample standard deviation tests
whether the returned uncertainty has the appropriate scale.  At nonzero injection,
the pull of Eq.~\eqref{eq:v4p9-pull} tests centering about the known amplitude and the
paired response of Eq.~\eqref{eq:v4p9-paired-response} separates recovery of the
added signal from a background-specific zero-signal offset.  Strengths
$z=0,1,3,5$ reuse the same background pseudo-spectrum and are consequently
correlated; they are one response curve, not four independent confirmations.

For each 100-background cell we report the usable count $n$, the pull mean and its
two-sided 90\% Student-$t$ interval, and the pull width with its exact normal-theory
90\% chi-square interval.  A cell is reported only if at least 95 of its 100 fits pass
the numerical criteria.  We account for multiplicity in two ways: a Holm correction
over the two planned replacement points at $z=0$, and a separate Holm correction over
the twenty displayed four-lane, five-mass points at $z=0$ after substitution.  The
correlated nonzero injections are response diagnostics and do not enter either family.
We flag a mean only when its Holm-adjusted two-sided $p$ value is below 0.05 and
$|\overline p|\geq0.2$.  The absence of such a flag is not proof that the bias is zero.

A practical centering tolerance, $|\overline p|<0.5$, was chosen after the
100-background results were available.  It is therefore a descriptive criterion, not
a predeclared equivalence test.  We report it alongside the $t$/Holm results, which are
more sensitive to small but statistically resolved offsets.  For the two replacement
points we also ask whether the full 90\% interval on the mean lies inside the practical
band.  A pull width is called consistent with one only when its 90\% interval contains
one.  None of these tests measures \CLs{} coverage.

\subsection{Training exclusion and blind-window choice}
\label{subsec:v491-guard-band-construction}

The GP must not be trained on the core or near tails of the signal it is later asked
to test.  A Gaussian signal retains non-negligible weight outside a
$1.64\sigma_m$ central interval, so training immediately beside that interval can
offer signal-like counts to the background fit and reduce the recovered amplitude.
The training-exclusion study therefore varied the sideband boundary before any
observed limit was used to choose it.

\begin{figure}[htbp]
\centering
\graphicorplaceholder{0.90\linewidth}{guard_band_figs/guard_band_5sigma_fake_data_gp_band_no_guard_2p25_extended2sigma.png}
\caption{Selected $2.25\sigma_m$ blind and training-exclusion geometry.  The GP is
trained outside the widened interval, while all bins inside it enter the local
signal-plus-background likelihood.  This retains the signal tails for extraction
without allowing those same bins to influence the sideband fit.}
\label{fig:v491-blind225-cartoon}
\end{figure}

\begin{table}[htbp]
\centering
\small
\begin{tabular}{ccccc}
\toprule
$k_{\rm train}$ & median expected $\eps^2_{95}$ proxy &
$\langle\widehat A/A_{\rm inj}\rangle_{z=3}$ &
$\langle\widehat A/A_{\rm inj}\rangle_{z=5}$ & max $f_{\rm train}$ \\
\midrule
1.96 & \num{7.92e-5} & 0.730 & 0.715 & 0.0530 \\
2.25 & \num{9.53e-5} & 0.961 & 0.943 & 0.0258 \\
2.58 & \num{1.08e-4} & 1.052 & 1.036 & 0.0106 \\
3.00 & \num{1.19e-4} & 1.049 & 1.037 & 0.00288 \\
\bottomrule
\end{tabular}
\caption{Evidence used to choose the training exclusion.  The train-$2.25$ row
substantially reduces the fraction of the signal template available to the GP while
recovering the injected amplitude and retaining useful local precision.  The first
column is a Gaussian-equivalent comparison proxy, not the final CLs limit.  The
complete note retains the mass-resolved reach, leakage, recovery, and pull batteries.}
\label{tab:v491-guard-band-choice-summary}
\end{table}

The widened $2.25\sigma_m$ interval was then used for both the training exclusion
and the extraction likelihood.  Compared with a $1.64\sigma_m$ extraction plus a
separate $1.64$--$2.25\sigma_m$ guard band, this construction gives comparable
closure while retaining the near signal tails in the likelihood.  The production
outcome is a single, explicit mass-dependent geometry for the background and signal
fits.

\subsection{Generating spectra and the role of functional forms}

The principal background toys are generated independently of the GP used for their
extraction.  A positive analytic intensity is fitted to a native HPS mass spectrum,
normalized on the specified toy support, integrated over every histogram bin, and
Poisson fluctuated.  A complete signal template is then added when $z>0$.  Every
pseudo-spectrum is subsequently analyzed with a new GP sideband fit.  This analytic
truth/full-refit construction is the central background-model test because success is
not guaranteed by generating toys from the GP model that will later fit them.

The earlier 2015 and 2016 studies used independently fitted shifted
sigmoid--power--exponential families.  The initial 2021 catalogue selected an
unshifted sigmoid--power--exponential model (internal label fSigPowExpQ) for the native
1\% source.  The later four-lane study used a generalized-gamma threshold model
(internal label fGenGammaThresh), fitted independently to the native 1\% and native
10\% sources.  These choices span smooth threshold, power-law, and broad-tail
curvature without adding resonance-scale structure.  Figure~\ref{fig:v491-functional-form-seed-overview}
shows the source spectra, selected means, and Poisson-ensemble construction.

\begin{figure}[p]
\centering
\graphicorplaceholder{0.99\linewidth}{toy_generation_figs/funcform_publication_overview.png}
\caption{Analytic generating means used in the year-matched functional-form studies.
Each dataset row compares the source spectrum with the candidate families and then
shows the selected mean against the generated toy ensemble.  The functions define
smooth positive pseudo-data means and are not asserted to be unique physical models
of the continuum.}
\label{fig:v491-functional-form-seed-overview}
\end{figure}

The mass-resolved injection scans selected the lower RBF length-scale factors without
using the observed limit curve.  For the frozen settings, the aggregate zero-signal
mean pulls are $+0.029$, $-0.016$, and $-0.005$ for the 2015, 2016 10\%, and 2021
1\% studies, respectively; the corresponding nonzero-injection signed averages are
$-0.223$, $-0.262$, and $-0.034$.  Mass-resolved mean, width,
significance-residual, and recovery checks support these summaries.  They motivate the
chosen lower factors but do not replace direct mass-by-mass
closure.

\subsection{Four-lane 2021 full-refit ensemble}
\label{sec:v491-four-lane-ensemble}

The full-100 2021 validation uses four exposure lanes and five mass hypotheses.  The
historical baseline is the frozen v4.2/v4.5 matched-refit contract: 40--300~MeV GP
support, 50--250~MeV search range, $k_{\min}=1.1$, $k_{\max}=15$,
five-bin local rebinning, log-mass/log-count preprocessing, the
$2.25\sigma_m$ blind and training exclusions, two required sidebands, and twelve
randomized optimizer restarts in addition to the nominal start.  The original four
lanes use the common generalized-gamma threshold model fitted independently to the
native 1\% and native 10\% sources.  Their source Pearson
$\chi^2/{\rm ndf}$ values are 1.108 and 2.781, respectively.  The latter
misdescription is largest through the steep threshold
turn-on and motivates a separately qualified 65~MeV replacement rather than a
post-hoc change to all other masses.

\begin{table}[htbp]
\centering
\small
\begin{tabularx}{\textwidth}{P{0.17\textwidth}P{0.20\textwidth}P{0.24\textwidth}Y}
\toprule
Displayed cell & Generating mean & GP support & Reason for scheme \\
\midrule
$1\%\times10$, $m\neq65$~MeV & generalized-gamma threshold model & 40--300~MeV & historical 100-background functional-form suite \\
$1\%\times10$, 65~MeV & degree-five threshold-turn-on continuation & 40--300~MeV & dedicated near-threshold description extended to 100 backgrounds \\
$1\%\times100$, all masses & generalized-gamma threshold model & 40--300~MeV & historical 100-background functional-form suite \\
native 10\%, $m\neq65$~MeV & generalized-gamma threshold model & 40--300~MeV & historical 100-background functional-form suite \\
native 10\%, 65~MeV & threshold-focused sigmoid--power--exponential mean & 30--300~MeV & includes the measured 30--40~MeV turn-on omitted by the historical support \\
$10\%\times10$, all masses & generalized-gamma threshold model & 40--300~MeV & historical 100-background functional-form suite \\
\bottomrule
\end{tabularx}
\caption{Composition of the 2021 background-validation display.  The two
special 65~MeV cells replace the corresponding historical cells; they are never
pooled with them.  Distinct markers in Figs.~\ref{fig:v491-zero-signal-full} and
\ref{fig:v491-pull-summary} identify the change of generating truth and, for native
10\%, the change of GP support.}
\label{tab:v491-composite-scheme-map}
\end{table}

\subsection{Targeted validation of the rising threshold}
\label{sec:v491-threshold-validation}

\paragraph{$1\%\times10$ threshold-turn-on truth.}
The first replacement is the source-fitted model developed in the dedicated
near-threshold study.  Over 35--80~MeV it is a logistic turn-on multiplied by an exponentiated
degree-five Chebyshev expansion and is joined smoothly to the specified high-mass tail.
The degree is selected before any GP extraction: the native-bin deviance per degree of
freedom is 0.994 for the native 1\% source and the five-bin value is 0.850, with no
parameter-bound contact.  Its disclosed 80--300~MeV tail Poisson deviance per degree
of freedom is 1.122, so this selected 1\%-source truth is broadly data-like over the
complete 40--300~MeV support rather than only in the fitted turn-on interval.  The
earlier 20-background result at 65~MeV had mean pull
0.300 and width 1.109; version 4.9.1 replaces that pilot with a predeclared
100-background continuation and reports 90\% intervals.

\paragraph{Native-10\% 30--300~MeV support.}
The second replacement models the measured 30--40~MeV shoulder and analyzes it with
the complete 30--300~MeV GP support.  Its threshold component is a logistic turn-on
times a degree-six exponentiated Chebyshev form fitted over 30--80~MeV, joined with a
$C^2$ taper to the frozen smooth anchor.  In the native 10\% source, the local
native-bin and five-bin deviances per degree of freedom are 1.140 and 1.421.  Thus the
generating spectrum is data-like in the threshold-sensitive region even though the
simple historical family is not.  Its 85--300~MeV tail deviance per degree of freedom
is 6.220; this native-10\% lane therefore relies on the explicit weak-correlation and
tail-influence qualification below rather than on a support-wide goodness-of-fit
claim.

This study tests the 65~MeV extraction under the 30--300~MeV support considered in
v4.9.1; it does not determine the 36--300~MeV support adopted in v4.9.5.  The question
here is whether a data-like generating mean that passes the local
source-fit tests can be processed by the complete extraction without material signal
bias.  With log preprocessing, the direct RBF correlation between masses $m$ and
$m'$ is
\begin{equation}
 \rho(m,m')=\exp\!\left[-\frac{\log^2(m'/m)}{2\ell^2}\right].
\end{equation}
At 65~MeV, the accepted optimized length scales in the two actual replacement
ensembles span 0.2784--0.3185 for $1\%\times10$ and 0.2468--0.2691 for native
10\%, with medians 0.2963 and 0.2592.  The corresponding direct correlations with
150~MeV span 1.10--3.19\% and 0.321--0.800\%, respectively.  Across the complete
55--70~MeV set the largest direct correlations with 150~MeV occur at 70~MeV and are
5.71\% and 1.81\%.  The correlation decreases rapidly above 150~MeV.  These results quantify why a
threshold-focused source fit can provide a meaningful local test while the generated
spectrum and extraction still retain the full GP support.

The correlation calculation is a fixed-hyperparameter statement, not a claim that
high-mass bins have mathematically zero influence.  All retained bins still enter the
global marginal likelihood and can shift the optimized length scale or kernel
constant.  A same-input diagnostic makes the aggregate effect explicit: after fixing
the selected hyperparameters, removing the training bins above 150~MeV shifts the
65~MeV pull by a mean of $-0.130$ with a paired 90\% interval
$[-0.151,-0.110]$ in $1\%\times10$, and by $-0.195$
$[-0.211,-0.179]$ in native 10\%.  The maximum absolute shifts are 0.376 and
0.400.  In the first lane the mean changes from 0.139 to 0.009; in the second it
changes from $-0.246$ to $-0.441$, whose 90\% interval
$[-0.611,-0.271]$ crosses the adopted practical boundary.  Counts above 150~MeV
therefore have modest, statistically resolved collective influence even though every
pairwise RBF correlation is small.  This diagnostic holds the chosen hyperparameters
fixed and does not test their global-likelihood response to tail removal.

The optimized kernels nevertheless explain why a threshold-focused source fit is
scientifically informative for the complete card: the low-mass extraction is
directly coupled most strongly to the turn-on region, while the generated spectrum,
specified smooth continuation, and full extraction retain the actual support.  The
restricted source fit thus provides a targeted conditional validation of local
65~MeV behavior under that full-card construction.  It does not establish
tail-independence or universal GP closure for every possible high-mass truth.

\subsection{Matched refits and numerical selection}

For every background, mass, and scheme, an independently optimized background-only
fit supplies the matched reference uncertainty.  Signals with
$z=1,3,5$ are generated with mean amplitude
$A_{\rm inj}=z\sigma_{A,{\rm ref}}$, their full Gaussian templates are added to the
same background counts, and the GP is refitted after applying the same training mask.
The injected amplitude is therefore calibrated to the local background-only fit,
whereas the pull denominator is always the uncertainty returned by the injected fit.

Whether a fit is kept never depends on its fitted signal or pull.  Repeated attempts
use the same pseudo-data with independent optimizer starts.  We choose among them using
the GP log marginal likelihood, reproducibility of the length scale and kernel
constant, covariance validity, uncertainty-ratio checks, and contact with the kernel
bounds.  The choice never uses $\widehat A$, the pull, agreement with the injected
value, $\eps^2$, or a local $p$ value.  A statistical outlier is retained whenever it
passes these numerical checks.

The historical study attempted 8000 fits and retained 7994; every cell contains 99 or
100 usable fits.  The six rejected fits failed reproducibility checks.  Each new
continuation contains 400 usable fits---100 backgrounds at each of $z=0,1,3,5$---with
no rejection and no kernel-bound contact.  Reaching those solutions required 614
optimizer attempts in the $1\%\times10$ threshold-turn-on lane and 610 in the native
10\% lane.  Failed starts remain recorded, and checks of the covariance, counts, and
pull calculation agree.  Numerical selection is therefore independent of the physics
result.

\subsection{Full-100 spurious-signal results}
\label{sec:v491-spurious-signal}

Table~\ref{tab:v491-sixtyfive-results} gives the two targeted 65~MeV results.  The
historical rows are shown only to explain why each replacement was required; no old
and new pulls are combined.  Both replacement samples exceed the 95-of-100 reporting
threshold and pass the numerical checks above.  We compare the results under the
original statistical test, the post-result practical
$|\overline p|<0.5$ criterion, and the separate unit-width diagnostic.

\begin{table}[htbp]
\centering
\small
\begin{tabularx}{\textwidth}{P{0.27\textwidth}cP{0.20\textwidth}P{0.20\textwidth}Y}
\toprule
Lane and 65~MeV scheme & $n$ & $\overline p$ [90\% CI] & $s_p$ [90\% CI] & mean assessment \\
\midrule
$1\%\times10$, historical & 100 & 0.7890 [0.6293,0.9488] & 0.9622 [0.8625,1.0907] & practical fail \\
$1\%\times10$, threshold turn-on, full & 100 & 0.1394 [$-0.0499$,0.3288] & 1.1404 [1.0222,1.2927] & practical pass \\
$1\%\times10$, threshold turn-on, reserve & 80 & 0.0993 [$-0.1150$,0.3135] & 1.1514 [1.0196,1.3265] & practical pass \\
native 10\%, historical & 100 & 0.7160 [0.5452,0.8867] & 1.0284 [0.9218,1.1657] & practical fail \\
native 10\%, 30--300 support, full & 100 & $-0.2463$ [$-0.4165$,$-0.0761$] & 1.0253 [0.9190,1.1622] & practical pass$^\dagger$ \\
native 10\%, 30--300 support, reserve & 75 & $-0.2840$ [$-0.4943$,$-0.0737$] & 1.0934 [0.9646,1.2661] & practical pass \\
\bottomrule
\end{tabularx}
\caption{Targeted 65~MeV background-only closure.  All headline intervals are
two-sided 90\% intervals: Student-$t$ for the mean and exact normal-theory
chi-square for the population standard deviation.  Historical entries are unpaired
context and are replaced, not pooled, in the composite validation figures.  Both
replacement full-100 mean intervals lie entirely within the post-result
$[-0.5,0.5]$ practical band, whereas the two historical means do not.
$^\dagger$The native-10\% full sample nevertheless has targeted Holm-2
$p=0.03632$ and $|\overline p|>0.2$; statistical resolution of a small offset and
practical centering are reported as distinct statements.  The $1\%\times10$ width
interval excludes one in both the full and reserve samples.}
\label{tab:v491-sixtyfive-results}
\end{table}

The 100-background continuations sharpen, rather than simply confirm, the pilot
interpretation.  The $1\%\times10$ threshold-turn-on truth is practically centered
in the full sample and in the independent reserve-80, but its pull distribution is
mildly overdispersed: the 90\% width interval excludes one in both samples.  Native
10\% with 30--300~MeV support moves from the historical $+0.716$ response to
$-0.246$; its full and reserve mean intervals both remain inside the practical band.
The native full sample is statistically displaced under the targeted two-endpoint
$t$/Holm screen, and the reserve-75 independently reproduces the negative direction.
This is evidence for a small signed offset, not a failure of the adopted practical
mean-closure margin.  Moving the support and improving the turn-on description reduces
the original 65~MeV offset enough to meet the post-result practical criterion, while
exact unit pull-width calibration remains a separate caveat.

After substitution, all twenty displayed zero-signal point means satisfy
$|\overline p|<0.5$; the largest absolute value is 0.338 at the historical
$10\%\times10$, 180~MeV cell.  Under the separate Holm-20 inferential screen that
180~MeV cell remains the only flag, while the native-10\%, 65~MeV composite adjusted
value is 0.291.  The larger multiplicity family is not used to erase the targeted
Holm-2 result.  Rather, the note distinguishes practical magnitude from evidence that
a nonzero offset is statistically resolved.  The original two 65~MeV cells failed the
practical margin with means 0.789 and 0.716; their replacements do not.

The same practical pattern holds when the correlated nonzero strengths are displayed:
all 80 lane--mass--strength point means lie inside $[-0.5,0.5]$, and 75 of their 80
pointwise 90\% intervals are wholly contained.  Exact unit-width calibration is less
uniform.  Three of the twenty zero-signal width intervals and fifteen of all eighty
width intervals exclude one.  At zero injection these are the replacement
$1\%\times10$, 65~MeV cell, historical $1\%\times10$ at 90~MeV, and historical
native 10\% at 120~MeV.  Thus practical mean/background-model-bias closure and exact
normal-pull calibration are separate outcomes of the same figure.

\ifwritingsample
\else
To preserve established cross-references in the complete analysis note, the three
consolidated displays below retain their earlier figure numbers.
\setcounter{figure}{63}
\fi
\begin{figure}[p]
\centering
\graphicorplaceholder{0.99\linewidth}{toy_generation_figs/v4p9p1_2021_background_validation_20260817/figure64_spurious_signal_consolidated_full100_90cl.pdf}
\caption{Full-100 zero-signal diagnostics for the four 2021 exposure lanes.  The
historical circular markers use the generalized-gamma threshold ensembles,
fitted independently to the native 1\% and native 10\% sources.
At 65~MeV only, distinct diamond and star markers replace---and are not pooled with---the
historical values: $1\%\times10$ uses the full-100 dedicated threshold-turn-on
continuation of the dedicated near-threshold study, while native 10\% uses the
full-100 threshold-focused extraction with 30--300~MeV GP support.  The substituted
points form a different, explicitly labeled validation scheme because the original
analytic family did not adequately model the rising threshold.  The panels retain
all accepted pulls, two-sided 90\% Student-$t$ intervals for means, and exact
normal-theory 90\% chi-square intervals for widths; no accepted pull is clipped.
The gray mean band is the post-result $\pm0.5$ practical centering tolerance.
Connecting lines are visual guides only.}
\label{fig:v491-zero-signal-full}
\end{figure}

\begin{figure}[p]
\centering
\graphicorplaceholder{0.99\linewidth}{toy_generation_figs/v4p9p1_2021_background_validation_20260817/figure65_onepctx10_65mev_table17_vs_historical_90cl.pdf}
\caption{Targeted $1\%\times10$, 65~MeV confirmation for the dedicated
threshold-turn-on truth compared with the historical generalized-gamma threshold
ensemble.  The upper panels show the zero-injection pull
distribution and the running mean; the lower panels show the mean and width across
the correlated $z=0,1,3,5$ extractions.  Mean intervals are two-sided 90\%
Student-$t$ intervals and width intervals are exact normal-theory 90\% chi-square
intervals.  The full-100 sample contains the original 20-background pilot; the
reserve-only result in Table~\ref{tab:v491-sixtyfive-results} demonstrates the same
practical centering and the same mild overdispersion.  Historical and replacement pulls are compared, not
pooled.  The gray $\pm0.2$ bands in this legacy-style focused display show the
inherited material-effect reference; they are not the post-result $\pm0.5$ practical
centering band used in this section.  The native-10\% replacement is shown in Figs.~\ref{fig:v491-zero-signal-full}
and \ref{fig:v491-pull-summary} and in Table~\ref{tab:v491-sixtyfive-results}.}
\label{fig:v491-sixtyfive-confirmation}
\end{figure}

\begin{figure}[p]
\centering
\graphicorplaceholder{0.99\linewidth}{toy_generation_figs/v4p9p1_2021_background_validation_20260817/pull_means_widths_consolidated_2x4_full100_90cl.pdf}
\caption{Compact composite background-validation summary.  Columns correspond to
$1\%\times10$, $1\%\times100$, native 10\%, and $10\%\times10$; the upper row
shows pull means and the lower row pull widths as functions of mass.  The four series
are $z=0,1,3,5$, with zero injection visually emphasized and the correlated nonzero
response series drawn more lightly.  Circular markers are the original full-100
common-family generalized-gamma threshold ensembles.  At 65~MeV, diamonds in
$1\%\times10$ and stars in native 10\% replace the complete $z=0,1,3,5$
historical sets with the targeted studies defined in
Table~\ref{tab:v491-composite-scheme-map}.  Mean bars are
two-sided 90\% Student-$t$ intervals and width bars are exact normal-theory two-sided
90\% chi-square intervals.  The substituted points are conditionally validated under
different source-fitted truths; they are not a homogeneous pooled ensemble, and the
four strengths share backgrounds.  Connecting lines are omitted across scheme
replacements.  The gray mean bands show the post-result $\pm0.5$ practical
centering tolerance.}
\label{fig:v491-pull-summary}
\end{figure}

Figure~\ref{fig:v491-sixtyfive-confirmation} directly compares the earlier and targeted
65~MeV ensembles, showing why the threshold-specific study was needed.

\subsection{Recovery of injected signals}

A centered zero-signal ensemble is necessary but not sufficient: a background model
could remove the false offset and still absorb a real narrow signal.  The nonzero
rows therefore repeat the full GP fit after adding matched one-, three-, and
five-reference-sigma signals.  Figure~\ref{fig:v491-pull-summary} displays the mean
and width at every strength and mass; Table~\ref{tab:v491-injection-recovery} gives
the numerical 65~MeV replacement results, including the paired response
$R=(\widehat A_z-\widehat A_0)/A_{\rm inj}$.

For the historical four-lane suite, the paired response separates signal recovery
from the additive background offset.  In the two replacements, the mean response is
stable at 0.966--0.969 over $z=1,3,5$, with mutually consistent 90\% intervals and no
worsening with strength.  The full procedure therefore recovers 96.6--96.9\% of the
tested amplitude, corresponding to a stable 3.1--3.4\% attenuation.  This is
disclosed as a conditional response calibration, not rounded to exact unity.
Agreement across strengths is not counted multiple times because the same 100
background realizations are reused.

\begin{table}[htbp]
\centering
\small
\begin{tabular}{llccc}
\toprule
65~MeV scheme & $z$ & $\overline p$ [90\% CI] & $s_p$ [90\% CI] & $\overline R$ [90\% CI] \\
\midrule
$1\%\times10$, threshold turn-on & 0 & $0.139$ [$-0.050$,0.329] & 1.140 [1.022,1.293] & --- \\
 & 1 & $0.107$ [$-0.082$,0.297] & 1.141 [1.023,1.293] & 0.9678 [0.9643,0.9713] \\
 & 3 & $0.039$ [$-0.149$,0.227] & 1.134 [1.016,1.285] & 0.9665 [0.9646,0.9684] \\
 & 5 & $-0.025$ [$-0.216$,0.165] & 1.147 [1.028,1.300] & 0.9670 [0.9654,0.9686] \\
\addlinespace
native 10\%, 30--300 support & 0 & $-0.246$ [$-0.417$,$-0.076$] & 1.025 [0.919,1.162] & --- \\
 & 1 & $-0.280$ [$-0.450$,$-0.109$] & 1.025 [0.919,1.162] & 0.9664 [0.9629,0.9700] \\
 & 3 & $-0.346$ [$-0.515$,$-0.177$] & 1.019 [0.913,1.155] & 0.9666 [0.9649,0.9684] \\
 & 5 & $-0.402$ [$-0.571$,$-0.234$] & 1.015 [0.909,1.150] & 0.9686 [0.9669,0.9703] \\
\bottomrule
\end{tabular}
\caption{Full-100 65~MeV extraction at all matched-reference strengths for the two
replacement schemes.  Mean intervals are two-sided 90\% Student-$t$ intervals and
width intervals are exact normal-theory two-sided 90\% chi-square intervals.  The
response summaries are paired background-by-background and use two-sided 90\%
Student-$t$ intervals.  All strengths share the same backgrounds and are correlated.  The
$1\%\times10$ widths are mildly high at every strength; the native-10\% widths are
compatible with one at this precision.}
\label{tab:v491-injection-recovery}
\end{table}

\subsection{What the validation establishes}

Taken together, these tests show practical mean closure and stable signal recovery for
the source-like spectra that were studied.  The training exclusion keeps the GP from
learning the signal tails, the pseudo-data means are fitted independently of the GP,
and every pseudo-spectrum receives a fresh sideband fit.  Numerical branches are
chosen without reference to the fitted signal or pull.  Near threshold, the historical
analytic family is replaced only in the two cells where its description of the source
is inadequate; tables and figures identify those replacements explicitly.

\begin{table}[htbp]
\centering
\small
\begin{tabularx}{\textwidth}{P{0.25\textwidth}P{0.22\textwidth}Y}
\toprule
Validation component & Evidence & Interpretation \\
\midrule
Training leakage and blind geometry & train-$2.25$, widened extraction window, recovery and leakage scans & supports the production mask \\
Functional-form source description & source-fit and bin-integral diagnostics; threshold-specific replacements at 65~MeV & the 1\%-source threshold model is data-like over the support; the native-10\% model is locally data-like and conditional on its tail \\
Spurious-signal means & Fig.~\ref{fig:v491-zero-signal-full} and Table~\ref{tab:v491-sixtyfive-results} & all 20 zero-signal point means meet $|\overline p|<0.5$; both replacement 90\% intervals lie inside the band \\
Pull widths & exact normal-theory 90\% chi-square intervals in Fig.~\ref{fig:v491-pull-summary} & exact unit-width calibration is not global: 3/20 zero-signal and 15/80 all-strength intervals exclude one \\
Signal injection & paired $z=1,3,5$ response with a fresh GP fit & stable 96.6--96.9\% response; tests extraction, not direct confidence-limit coverage \\
Numerical stability & multistart and reproducibility checks independent of the fitted signal & all 800 fits in the new lanes are usable \\
High-mass influence on threshold & direct RBF correlations and fixed-hyperparameter tail-removal diagnostic & local validation is meaningful, but the aggregate tail effect is nonzero and tail independence is not claimed \\
Observed inference & no observed spectrum or limit enters truth/branch selection & unchanged by these conditional toys \\
\bottomrule
\end{tabularx}
\caption{Summary of the background-validation evidence.  The practical mean
criterion, statistical tests, pull-width calibration, signal response, and tail
influence are kept separate; none is promoted into a coverage or universality claim.}
\label{tab:v491-background-validation-decision}
\end{table}

Under the adopted practical criterion, the near-threshold concern is resolved within
this source-conditioned suite.  The two replacement 100-background intervals and their
independent reserve samples lie inside the $\pm0.5$ mean-pull band.  So do all twenty
zero-signal point estimates and all eighty correlated injected-strength point estimates
in the four-lane display.  The paired response remains 96.6--96.9\%.  The practical
mean-bias requirement is therefore met for these tests, provided that each replacement
remains identified by its generating model and GP support.

The stricter diagnostics remain important.  The original $10\%\times10$, 180~MeV
cell retains an exploratory Holm-20 flag with mean 0.338, and the native-10\%, 65~MeV
replacement retains a targeted Holm-2 flag.  Three of twenty zero-signal width intervals
and fifteen of eighty all-strength width intervals exclude one; the stable signal
attenuation is 3.1--3.4\%; and the fixed-hyperparameter tail-removal shifts are nonzero.
The study therefore does not prove zero bias for arbitrary smooth spectra, exact pull
calibration, confidence-limit coverage, bias in the observed data, or a scan-global
significance calibration.  Any larger study should retain the source, support,
optimizer, and post-result tolerance labels used here.
